\clearpage
\section{Arithmetic and exact forward execution padding}
\label{sec:copying-arithmetic-padding}
All three arithmetic operations now have actual local graph rules and checked
executions in the combined machine. Addition and subtraction evaluate both
numeric arguments in order, perform one rewrite and return the numeric result
in seven transitions. Unary minus takes five transitions. This completes the
local rule families for the existing example alphabet; it does not prove
unique matching, determinism or refocusing of the complete copying UAM.

The same milestone proves forward padding of every combined transition and
every finite run, including their actual contextual matches. A run of length
$n$ gives a run of exactly length $n$ after an unused allowance layer is added.
Successful termination is preserved under arbitrarily many layers, and when
both a closing context and its inserted program are padded. This proves one
direction of the outstanding padding obligation. It does not reflect arbitrary
runs from a larger allowance, so the two-sided execution-padding interface
remains unproved. Neither direction has been inferred from determinism.

The terms of Proposition~12.2 and the Figure~16 translation are unchanged.
The natural-number arithmetic convention below is explicit; no new numeric
constant datatype or altered term translation is introduced.

\Needspace{12\baselineskip}
\subsection{Arithmetic meanings and literal local rules}

\termdef{natural-arithmetic-convention}
The existing constants carry natural numbers. Binary arithmetic codes select
addition or subtraction and map to the corresponding original operation label.
Their evaluation functions are $m+n$ and $m\mathbin{\dot{-}}n$, where
$\dot{-}$ denotes natural subtraction truncated at zero. Unary minus is
$0\mathbin{\dot{-}}n$. The project author clarified the intended equation
$-x=0-x$ during this milestone; keeping the existing natural-number domain
therefore makes its result zero for every $n$. This last fact is proved with
Lean's natural-subtraction theorem, rather than built into the syntax.
Figure~10(a) explicitly depicts binary arithmetic; the unary rule uses this
clarified interpretation of Figure~13's unary operation.

\lean{Copying/Arithmetic}{9}{11}{Arithmetic.BinaryCode}

\lean{Copying/Arithmetic}{13}{15}{Arithmetic.BinaryCode.symbol}

\lean{Copying/Arithmetic}{17}{19}{Arithmetic.BinaryCode.eval}

\lean{Copying/Arithmetic}{22}{22}{Arithmetic.neg}

\lean{Copying/Arithmetic}{24}{24}{Arithmetic.neg_eq_zero}

\termdef{numeric-arithmetic-redex}
A number is the original constant corolla, with one term input and no outputs.
The unary redex is a unary-minus corolla serially connected to that number.
The binary redex is the chosen arithmetic corolla connected to the ordered
tensor of two number corollas. Their reducts are single number corollas with
the computed result. Thus the binary construction is the three-edge to
one-edge rule of Figure~10(a). The rule relation has one unary constructor
and one constructor parameterized by the two binary codes. It does not accept
Boolean, unit, function or reference operands as numbers. Their separate
language operations and search behaviour remain present.

\lean{Copying/Arithmetic}{26}{26}{Arithmetic.number}

\lean{Copying/Arithmetic}{27}{28}{Arithmetic.unaryBefore}

\lean{Copying/Arithmetic}{29}{29}{Arithmetic.unaryAfter}

\lean{Copying/Arithmetic}{30}{31}{Arithmetic.binaryBefore}

\lean{Copying/Arithmetic}{32}{32}{Arithmetic.binaryAfter}

\lean{Copying/Arithmetic}{35}{37}{Arithmetic.Rule}

A padded number corolla is literally the same primitive corolla at the
next allowance. The established serial-padding isomorphism therefore gives
unary-redex padding. Tensor padding and serial padding give binary-redex
padding. Each result-padding equation is the number equation specialized to
the computed result. The isomorphism form of number padding is obtained by
rewriting that equality and taking the identity isomorphism. These facts
supply the new arithmetic cases of combined operational padding below.

\lean{Copying/Arithmetic}{39}{39}{Arithmetic.number_raise}

\lean{Copying/Arithmetic}{41}{44}{Arithmetic.unaryBefore_raise}

\lean{Copying/Arithmetic}{46}{46}{Arithmetic.unaryAfter_raise}

\lean{Copying/Arithmetic}{48}{51}{Arithmetic.binaryBefore_raise}

\lean{Copying/Arithmetic}{53}{54}{Arithmetic.binaryAfter_raise}

\lean{Copying/Arithmetic}{67}{70}{Arithmetic.number_raise_iso}

Searching a number corolla returns immediately at the same external
input. The proof supplies its actual edge, source membership and the primitive
entry equation. A number is a stable constant, whose established stability
consequences prove absence of focus and hole labels. This certifies its
safety independently of any termination equivalence.

\lean{Copying/Arithmetic}{57}{60}{Arithmetic.number_returns}

\lean{Copying/Arithmetic}{63}{64}{Arithmetic.safe_number}

\Needspace{12\baselineskip}
\subsection{Ordered search and concrete computation}

The unary head and numeric-value selectors name the two actual edges;
the argument selector names their internal connection. Entry searches that
connection. The numeric edge returns the token, and backtracking through the
unary head requests rewriting at the external input. Each proof checks
incidence and the relevant entry or resume function, separately for allowance
zero and positive allowances where the primitive carrier presentations differ.

\lean{Copying/ArithmeticExecution}{9}{9}{Arithmetic.unaryHead}

\lean{Copying/ArithmeticExecution}{10}{10}{Arithmetic.unaryValue}

\lean{Copying/ArithmeticExecution}{11}{11}{Arithmetic.unaryArgument}

\lean{Copying/ArithmeticExecution}{13}{17}{Arithmetic.unary_enter}

\lean{Copying/ArithmeticExecution}{19}{23}{Arithmetic.unary_value_returns}

\lean{Copying/ArithmeticExecution}{25}{30}{Arithmetic.unary_ready}

The binary selectors name the operation head, either operand edge
and either ordered internal argument coordinate. The five search steps are:
enter the left argument, return from its number, resume into the right
argument, return from its number, and request rewriting at the root. The
argument indices are zero and one, and their bounds are proved from the
operation's arity. These are ordered incidence proofs, not an evaluation
oracle.

\lean{Copying/ArithmeticExecution}{32}{33}{Arithmetic.binaryHead}

\lean{Copying/ArithmeticExecution}{34}{38}{Arithmetic.binaryValue}

\lean{Copying/ArithmeticExecution}{39}{40}{Arithmetic.binaryArgument}

\lean{Copying/ArithmeticExecution}{42}{46}{Arithmetic.binary_enter}

\lean{Copying/ArithmeticExecution}{48}{52}{Arithmetic.binary_left_returns}

\lean{Copying/ArithmeticExecution}{54}{59}{Arithmetic.binary_search_right}

\lean{Copying/ArithmeticExecution}{61}{65}{Arithmetic.binary_right_returns}

\lean{Copying/ArithmeticExecution}{67}{72}{Arithmetic.binary_ready}

Inversion of an arithmetic rule chooses its actual combined local
constructor. The existing local-match constructor then supplies a combined
rewrite from rewrite focus to search focus on the result. Appending that
transition and the one-step number return to the three or five search steps
gives the five- and seven-step executions. No standalone arithmetic execution
parameter is assumed. Since these rules are in the original combined
relation, its already-proved contextual lifting, recursive-isomorphism
transport, finite-run representation and token-typing theorems now cover
arithmetic too.

\lean{Copying/ArithmeticExecution}{75}{79}{Arithmetic.rule_step}

\lean{Copying/ArithmeticExecution}{82}{87}{Arithmetic.unary_run}

\lean{Copying/ArithmeticExecution}{90}{96}{Arithmetic.binary_run}

The input graphs satisfy the translation language's label predicate:
every primitive is an operation, and serial and tensor composition preserve
that predicate. Consequently the established language safety theorems yield
focus-free and hole-free certificates. Evaluating the actual finite edge
enumerations proves that a binary redex has three edges and its result one,
while a unary redex has two edges and its result one. The proofs hold for
all allowances and all natural operands.

\lean{Copying/ArithmeticExecution}{98}{100}{Arithmetic.unary_language}

\lean{Copying/ArithmeticExecution}{102}{105}{Arithmetic.binary_language}

\lean{Copying/ArithmeticExecution}{108}{110}{Arithmetic.binary_edge_counts}

\lean{Copying/ArithmeticExecution}{112}{114}{Arithmetic.unary_edge_counts}

\Needspace{12\baselineskip}
\subsection{Padding contextual attachments without moving the token}

\termdef{padding-attachment-law}
Uniform padding keeps each outer vertex and edge carrier and relabels its
edges, recursively padding box bodies. For a simultaneously expanded host,
the padded filling of an edge is the padded component with interface equality
casts supplied by the label-padding laws. Those casts are necessary because
the hole signature is indexed by the host edge's label. They change no port
index or wire type.

\lean{Copying/PaddingInteractions}{10}{14}{Padding.expandFilling}

The expansion attachment equation says that inserting the appropriately
cast component port into the padded host reaches exactly the original vertex.
For an internal port both sides are the same tagged component vertex. For
an input or output, both sides use the same host source or target position.
Case analysis on ports and the incidence-index definitions proves the equation.
The expansion-padding isomorphism uses identity vertex and edge bijections;
labels agree, and its source and target equations follow by composing list
maps and applying that attachment equation to each port.

\lean{Copying/PaddingInteractions}{17}{24}{Padding.expand_insert}

\lean{Copying/PaddingInteractions}{27}{46}{Padding.expand_graph}

Single-slot replacement has the analogous attachment equation. The
host embedding is unchanged and a cast inserted input or output reaches the
same host boundary incidence. The corresponding padding isomorphism again
uses identity outer bijections. Host-edge incidence commutes with the host
embedding; inserted-edge incidence commutes by the attachment equation.
The source filler requires no gratuitous boundary cast, which lets this
isomorphism compare the actual graphs used by the rewrite constructor.

\lean{Copying/PaddingInteractions}{49}{57}{Padding.plug_insert}

\lean{Copying/PaddingInteractions}{59}{85}{Padding.plug_graph}

Uniform padding also supplies a concrete search-transport structure.
Both port and edge bijections are identities. Case analysis on the allowance
and original label proves that an edge has a specified primitive label before
padding exactly when it has that same primitive label afterwards. Source and
target lists are unchanged, and external-input status is unchanged. The
previous search-transport theorem can therefore be reused without modifying
intrinsic search.

\lean{Copying/PaddingInteractions}{88}{97}{Padding.search}

\Needspace{12\baselineskip}
\subsection{Every combined rewrite and run pads forward}

First, replacing ordered interfaces by equal lists transports a rewrite
and casts its active coordinates. After eliminating the interface equalities,
the proof constructs identity-carrier isomorphisms to the reboundaried graphs
and invokes the existing representation closure. A separate helper observes
that graph isomorphisms fix each indexed external input. Another helper
expresses a previously proved rule-side padding isomorphism using uniform
padding instead of the equivalent original raising function.

\lean{Copying/CombinedPadding}{10}{26}{Combined.rewrite_reboundary}

\lean{Copying/CombinedPadding}{29}{31}{Combined.rewrite_input_iso}

\lean{Copying/CombinedPadding}{34}{36}{Combined.graphIsoOfRaise}

Every local rule pads forward. For allocation, assignment, micro-beta
and copying, the proof uses their actual redex and reduct padding
isomorphisms. Stability transports by the proved stability-padding equivalence;
copyability transports by its proved padding theorem. Dereference and the two
atom-equality cases agree directly with their next-allowance primitive graphs.
The unary and binary arithmetic cases use the new redex and number padding
laws. All active ports are external term inputs at this local stage, so their
coordinates agree by the input-isomorphism helper. No side condition is
dropped or replaced by a machine-determinism assumption.

\lean{Copying/CombinedPadding}{39}{69}{Combined.local_padding}

A recursive graph isomorphism lifts through uniform padding with its
original outer bijections. The recursive label-equivalence padding theorem
supplies its only new field; all incidence and internal-label laws are
unchanged. Induction on a contextual rewrite now handles every constructor.
The local case is the theorem above. The renaming case uses padded
isomorphisms. For replacement, cast the padded component's interfaces,
rewrite in the padded host, and rename through the replacement-padding
isomorphisms; the attachment equations identify the actual token coordinates.
For simultaneous expansion, all unselected components remain equal after
padding, and the corresponding expansion isomorphisms finish the same argument.
This establishes forward transport for arbitrary finite shallow matches.

\lean{Copying/CombinedPadding}{72}{79}{Combined.graphIso}

\lean{Copying/CombinedPadding}{82}{108}{Combined.rewrite_padding}

\termdef{forward-execution-padding}
Padding a state increases its allowance and pads its actual graph, while
retaining its mode and port. Finality is therefore equivalent before and
after this state transformation. A combined step lifts by cases: intrinsic
search uses the uniform search transport, and a graph-changing step uses
rewrite padding. Induction lifts every run to its literal padded endpoints,
with unchanged length. A successful run retains its final endpoint, giving
preservation of successful termination at each exact cost.

\lean{Copying/CombinedPadding}{111}{111}{Combined.padState}

\lean{Copying/CombinedPadding}{113}{114}{Combined.padState_final}

\lean{Copying/CombinedPadding}{117}{123}{Combined.step_padding}

\lean{Copying/CombinedPadding}{126}{130}{Combined.run_padding}

\lean{Copying/CombinedPadding}{132}{135}{Combined.terminatesIn_padding}

\Needspace{12\baselineskip}
\subsection{Unbounded programs and padded observations}

The unbounded-state padding function adds one to its allowance and
applies the fixed-allowance state transformation. Finality remains equivalent.
Inverting an unbounded step and applying fixed-allowance step padding gives
one lifted transition. Run induction gives the same number of unbounded
transitions. For any cost preorder $Q$, successful source runs therefore
witness refinement to the padded state using $Q(n,n)$. This is a forward
refinement theorem, not a two-sided equivalence.

\lean{Copying/CombinedPaddingExecution}{8}{8}{Combined.padUnbounded}

\lean{Copying/CombinedPaddingExecution}{10}{11}{Combined.padUnbounded_final}

\lean{Copying/CombinedPaddingExecution}{13}{15}{Combined.unbounded_step_padding}

\lean{Copying/CombinedPaddingExecution}{17}{21}{Combined.unbounded_run_padding}

\lean{Copying/CombinedPaddingExecution}{24}{27}{Combined.padding_refines}

Padding an initialized closed program agrees with initializing its
original padded program object. The proof applies equality congruence to
the graph-padding equality, with the search input fixed explicitly. Induction
on the number of added layers composes one-step refinements; no additional
machine transitions are introduced by this proof. Successful termination
then follows for every finite padding by applying the exact-cost refinement
to the original successful run.

\lean{Copying/CombinedPaddingExecution}{29}{31}{Combined.padding_initial}

\lean{Copying/CombinedPaddingExecution}{34}{41}{Combined.initial_pad_refines}

\lean{Copying/CombinedPaddingExecution}{44}{48}{Combined.terminates_pad}

Initial states of recursively isomorphic closed graphs have the same
search mode and external input, so the previously proved representation
relation gives equivalence for every cost preorder, including exact cost.
To pad a closed observation, first pad the already closed program using the
new forward refinement. The established context-filling padding isomorphism
then identifies that padded program with the result of closing the padded
fragment in the padded context. Composing with exact isomorphism invariance
proves forward preservation of all successful observation costs.

\lean{Copying/CombinedPaddingExecution}{51}{54}{Combined.initial_iso}

\lean{Copying/CombinedPaddingExecution}{57}{63}{Combined.close_pad_refines}

The converse needs an additional proof: an arbitrary run from a padded
graph may use a match presented at the larger allowance, and that match must
be reflected to the smaller representation. None of the forward constructors
provides this inversion automatically. Until reflection is proved, the
execution does not satisfy \texttt{RespectsPadding}; the cofinal observation
criterion and general semantic weakening cannot be concluded from this
milestone alone.

\Needspace{12\baselineskip}
\subsection{Arithmetic and representation regressions}

The unary and binary local redexes define actual closed-program objects,
using the language safety certificates proved above. Their finite runs lift
to the unbounded machine and give explicit five- and seven-transition
termination witnesses. These program objects denote the literal local
redex graphs. They are not advertised as the complete Figure~16 translations
of arithmetic expressions with surrounding environment distributors.

\lean{Copying/ArithmeticExamples}{11}{13}{Arithmetic.Examples.unaryProgram}

\lean{Copying/ArithmeticExamples}{15}{17}{Arithmetic.Examples.binaryProgram}

\lean{Copying/ArithmeticExamples}{19}{20}{Arithmetic.Examples.unary_terminates}

\lean{Copying/ArithmeticExamples}{22}{23}{Arithmetic.Examples.binary_terminates}

The unary-result theorem substitutes the proved zero result into its
actual endpoint state. The subtraction-underflow theorem does the same for
every pair $m\leq n$, using natural subtraction. Endpoint-state congruence is
used because the state contains a graph and a vertex on that graph. Addition
and subtraction with a zero right operand both return $n$ in seven steps.
All these regressions quantify over arbitrary natural numbers and allowances.

\lean{Copying/ArithmeticExamples}{26}{33}{Arithmetic.Examples.unary_result}

\lean{Copying/ArithmeticExamples}{36}{43}{Arithmetic.Examples.subtraction_underflow}

\lean{Copying/ArithmeticExamples}{45}{48}{Arithmetic.Examples.addition_zero}

\lean{Copying/ArithmeticExamples}{50}{53}{Arithmetic.Examples.subtraction_zero}

A contextual regression connects an outer unary-minus corolla to a
binary arithmetic redex. Its argument selector is the actual right-component
insertion coordinate. The outer head enters that argument in one search step.
The right-serial run-lifting theorem executes all seven inner transitions,
including its graph-changing arithmetic rule. Backtracking then resumes the
outer head, which performs its own rewrite and numeric return. The resulting
trace has $1+7+1+1+1=11$ transitions. This checks interaction of two real
rewrites across a contextual boundary, not merely an isolated compute step.

\lean{Copying/ArithmeticExamples}{56}{57}{Arithmetic.Examples.nested}

\lean{Copying/ArithmeticExamples}{59}{60}{Arithmetic.Examples.innerInput}

\lean{Copying/ArithmeticExamples}{62}{66}{Arithmetic.Examples.nested_enter}

\lean{Copying/ArithmeticExamples}{69}{76}{Arithmetic.Examples.nested_run}

The original translated reference program still terminates after any
finite padding, by the new general termination theorem. The application
regression with a literal translated constant body retains a successful run
of exactly seven transitions after any padding. Every binary arithmetic redex
also retains its exact seven-transition cost at arbitrarily higher allowances.
These proofs use the exact-cost refinement directly and eliminate its equality
of run lengths; they do not claim uniqueness of all combined executions.

\lean{Copying/ArithmeticExamples}{79}{81}{Arithmetic.Examples.translated_reference_pad}

\lean{Copying/ArithmeticExamples}{84}{92}{Arithmetic.Examples.application_pad_cost}

\lean{Copying/ArithmeticExamples}{95}{103}{Arithmetic.Examples.binary_pad_cost}

The full build checks 355 jobs. The three audits pass for 6 abstract,
477 graph and 912 copying declarations, with only standard Lean logical
axioms. Appendix~B.1 removes arithmetic construction and forward execution
padding from the unfinished obligations, retaining reflection of padded
matches/runs, matching independence, full-machine structural laws, template
robustness and final assembly. Estimated progress by mathematical effort is
about 75\%, broadly 60--85\% (26 October 2026). This is a judgement about proof
difficulty; the remaining template and machine arguments still carry substantial
uncertainty.

